\clearpage
\begin{examplebox}
\textbf{例 1：由切点直接写切线}

已知双曲线 $x^2/4-y^2/3=1$，求点 $P(4,3)$ 处的切线。

\textbf{识别：}已知的是切点，可以先验点，再用半代换。
\[
\frac{4^2}{4}-\frac{3^2}{3}=1,\qquad
\frac{4x}{4}-\frac{3y}{3}=1.
\]
故切线为 $\boxed{x-y=1}$，几何位置见首页图。

\textbf{三法互证：}（最低门槛）相减并让另一点靠近 $P$，切线斜率趋向
$3\cdot4/(4\cdot3)=1$；（最快）导数法同样得到 $k=1$；
（常规）令 $y=x+t$，由切点 $P(4,3)$ 得 $t=-1$，
代入双曲线得到下面的二重根等式。三种方法得到同一条直线。

\textbf{回验：}代入 $y=x-1$，得
\[
\frac{x^2}{4}-\frac{(x-1)^2}{3}=1
\quad\Longleftrightarrow\quad (x-4)^2=0.
\]
重根 $x=4$ 对应 $y=3$，与切点一致。

\medskip
\textbf{例 2：焦点在 $y$ 轴，保留正负项}

已知双曲线 $y^2/5-x^2/4=1$，求点 $P(2,\sqrt{10})$ 处的切线。

\textbf{识别：}正项是 $y^2/5$，因此半代换后也须以 $y_0y/5$ 为正项。
先验点：$10/5-4/4=1$；再代换：
\[
\frac{\sqrt{10}\,y}{5}-\frac{2x}{4}=1
\quad\Longleftrightarrow\quad
\boxed{2\sqrt{10}\,y-5x=10}.
\]
\textbf{回验：}代入切点，等式左边为 $20-10=10$。
曲线求导得 $y'=5x/(4y)$，在切点处为 $5/(2\sqrt{10})$，
也等于所得直线的斜率。
\end{examplebox}
\clearpage
\begin{examplebox}
\textbf{例 3：已知切线经过的点，反求切点}

双曲线 $x^2/2-y^2=1$ 的切线经过 $A(1,2)$，求所有切点。
\begin{center}
\textbf{过 $A$ 的两条切线，分别接触左右两支}\par\medskip
\begin{tikzpicture}[x=1.3cm,y=1.3cm,>=Stealth,font=\small]
\draw[->,gray] (-3.6,0)--(3.6,0) node[right] {$x$};
\draw[->,gray] (0,-2.7)--(0,3.3) node[above] {$y$};
\draw[blue,very thick,domain=-1.55:1.55,samples=100] plot ({sqrt(2)*cosh(\x)},{sinh(\x)});
\draw[blue,very thick,domain=-1.55:1.55,samples=100] plot ({-sqrt(2)*cosh(\x)},{sinh(\x)});
\draw[orange!85!black,thick] (-3.4,-2.4)--(2.2,3.2);
\draw[orange!85!black,thick] (0.76,3.2)--(1.9,-2.5);
\fill (1,2) circle (1.6pt);
\node at (-1,1.4) {$A(1,2)$};
\draw[gray,thin] (-.3,1.5)--(.94,1.97);
\fill (-2,-1) circle (1.6pt) node[below right=5pt] {$P_1(-2,-1)$};
\fill (1.428571,-.142857) circle (1.6pt);
\node[anchor=west] at (2.4,-.55) {$P_2(\frac{10}{7},-\frac17)$};
\draw[gray,thin] (1.5,-.17)--(2.3,-.5);
\node[blue,fill=BGExample,inner sep=3pt] at (-2.1,2.6) {$\dfrac{x^2}{2}-y^2=1$};
\end{tikzpicture}

\end{center}
\textbf{识别与列式：}$A$ 是切线经过的点，不能直接当作切点。
设切点为 $P(x_0,y_0)$，其切线为 $x_0x/2-y_0y=1$，故
\[
\begin{cases}
x_0^2/2-y_0^2=1, & \text{切点在曲线上},\\
x_0/2-2y_0=1, & \text{切线经过 }A.
\end{cases}
\]
\textbf{消元：}由 $x_0=2+4y_0$ 得
\[
\frac{(2+4y_0)^2}{2}-y_0^2=1
\ \Longrightarrow\ 7y_0^2+8y_0+1=(7y_0+1)(y_0+1)=0.
\]
所以 $\boxed{P_1=(-2,-1),\quad P_2=(10/7,-1/7)}$。

\textbf{回验：}两点均满足原曲线方程；对应切线为 $y=x+1$、$5x+y=7$，
均经过 $A(1,2)$。这两条切线与图中的两次接触一致。
\end{examplebox}
